\documentclass{amsart}
% For dealing with references we use the comment environment
\usepackage{verbatim}
\newenvironment{reference}{\comment}{\endcomment}
%\newenvironment{reference}{}{}
\newenvironment{slogan}{\comment}{\endcomment}
\newenvironment{history}{\comment}{\endcomment}

% For commutative diagrams we use Xy-pic
\usepackage[all]{xy}

% We use 2cell for 2-commutative diagrams.
\xyoption{2cell}
\UseAllTwocells

% We use multicol for the list of chapters between chapters
\usepackage{multicol}

% This is generally recommended for better output
\usepackage{lmodern}
\usepackage[T1]{fontenc}

% For cross-file-references

% Package for hypertext links:
\usepackage{hyperref}

% For any local file, say "hello.tex" you want to link to please

% Theorem environments.
%
\theoremstyle{plain}
\newtheorem{theorem}[subsection]{Theorem}
\newtheorem{proposition}[subsection]{Proposition}
\newtheorem{lemma}[subsection]{Lemma}

\theoremstyle{definition}
\newtheorem{definition}[subsection]{Definition}
\newtheorem{example}[subsection]{Example}
\newtheorem{exercise}[subsection]{Exercise}
\newtheorem{situation}[subsection]{Situation}

\theoremstyle{remark}
\newtheorem{remark}[subsection]{Remark}
\newtheorem{remarks}[subsection]{Remarks}

\numberwithin{equation}{subsection}

% Macros
%
\def\lim{\mathop{\mathrm{lim}}\nolimits}
\def\colim{\mathop{\mathrm{colim}}\nolimits}
\def\Spec{\mathop{\mathrm{Spec}}}
\def\Hom{\mathop{\mathrm{Hom}}\nolimits}
\def\Ext{\mathop{\mathrm{Ext}}\nolimits}
\def\SheafHom{\mathop{\mathcal{H}\!\mathit{om}}\nolimits}
\def\SheafExt{\mathop{\mathcal{E}\!\mathit{xt}}\nolimits}
\def\Sch{\mathit{Sch}}
\def\Mor{\mathop{\mathrm{Mor}}\nolimits}
\def\Ob{\mathop{\mathrm{Ob}}\nolimits}
\def\Sh{\mathop{\mathit{Sh}}\nolimits}
\def\NL{\mathop{N\!L}\nolimits}
\def\CH{\mathop{\mathrm{CH}}\nolimits}
\def\proetale{{pro\text{-}\acute{e}tale}}
\def\etale{{\acute{e}tale}}
\def\QCoh{\mathit{QCoh}}
\def\Ker{\mathop{\mathrm{Ker}}}
\def\Im{\mathop{\mathrm{Im}}}
\def\Coker{\mathop{\mathrm{Coker}}}
\def\Coim{\mathop{\mathrm{Coim}}}

% Boxtimes
%
\DeclareMathSymbol{\boxtimes}{\mathbin}{AMSa}{"02}

%
% Macros for moduli stacks/spaces
%
\def\QCohstack{\mathcal{QC}\!\mathit{oh}}
\def\Cohstack{\mathcal{C}\!\mathit{oh}}
\def\Spacesstack{\mathcal{S}\!\mathit{paces}}
\def\Quotfunctor{\mathrm{Quot}}
\def\Hilbfunctor{\mathrm{Hilb}}
\def\Curvesstack{\mathcal{C}\!\mathit{urves}}
\def\Polarizedstack{\mathcal{P}\!\mathit{olarized}}
\def\Complexesstack{\mathcal{C}\!\mathit{omplexes}}
% \Pic is the operator that assigns to X its picard group, usage \Pic(X)
% \Picardstack_{X/B} denotes the Picard stack of X over B
% \Picardfunctor_{X/B} denotes the Picard functor of X over B
\def\Pic{\mathop{\mathrm{Pic}}\nolimits}
\def\Picardstack{\mathcal{P}\!\mathit{ic}}
\def\Picardfunctor{\mathrm{Pic}}
\def\Deformationcategory{\mathcal{D}\!\mathit{ef}}

\setlength{\emergencystretch}{3em}
\begin{document}
\title[Stacks Project, Tag 04MK]{581 --- Reduced connected schemes locally of finite type over two fields have the same integral closure of their ground fields in global functions}
\maketitle
\noindent Copyright (C) 2005--2025 Johan de Jong. Stacks Project authors. Source: \href{https://stacks.math.columbia.edu/tag/04MK}{Tag 04MK}.

\noindent Editorial difficulty note (not author text). Difficulty H2: For quasi-compact schemes the proof uses finite generation of the unit group modulo constants and the Nullstellensatz to compare the two constant fields. The general case requires gluing local constant fields across connected overlaps and propagating a single monic relation through the connected covering graph..

\noindent Editorial provenance note (not author text). History: excerpt from revision \texttt{a0444\allowbreak 6e57e\allowbreak c1fbc\allowbreak 252a8\allowbreak 71afc\allowbreak ec775\allowbreak 2fb28\allowbreak 07b14}, varieties.tex, lines 5964--6056. Reference presentation was adapted; original-author context and core support blocks were retained or added. The mathematical wording is unchanged. Licensed under GNU FDL 1.2 or later; no invariant sections or cover texts. See ../licenses/COPYING. Context blocks reproduce author definitions and supporting results; they are not additional counted problems.

\begin{proposition}
\label{proposition-unique-base-field}
Let $X$ be a scheme. Let $a : X \to \Spec(k_1)$ and
$b : X \to \Spec(k_2)$ be morphisms from $X$ to spectra of fields.
Assume $a, b$ are locally of finite type, and
$X$ is reduced, and connected. Then we have
$k_1' = k_2'$, where $k_i' \subset \Gamma(X, \mathcal{O}_X)$ is
the integral closure of $k_i$ in $\Gamma(X, \mathcal{O}_X)$.
\end{proposition}

\begin{proof}
First, assume the lemma holds in case $X$ is quasi-compact (we will
do the quasi-compact case below).
As $X$ is locally of finite type over a field, it is locally Noetherian, see
Morphisms, Lemma \href{https://stacks.math.columbia.edu/tag/01T6}{lemma finite type noetherian (Tag 01T6)}.
In particular this means that it is locally connected,
connected components of open subsets are open, and
intersections of quasi-compact opens are quasi-compact, see
Properties, Lemma \href{https://stacks.math.columbia.edu/tag/01OZ}{lemma Noetherian topology (Tag 01OZ)},
Topology, Lemma \href{https://stacks.math.columbia.edu/tag/04ME}{lemma locally connected (Tag 04ME)},
Topology, Section \href{https://stacks.math.columbia.edu/tag/0050}{section noetherian (Tag 0050)}, and
Topology, Lemma \href{https://stacks.math.columbia.edu/tag/005L}{lemma constructible Noetherian space (Tag 005L)}.
Pick an open covering $X = \bigcup_{i \in I} U_i$
such that each $U_i$ is quasi-compact and connected.
For each $i$ let $K_i \subset \mathcal{O}_X(U_i)$ be the integral
closure of $k_1$ and of $k_2$.
For each pair $i, j \in I$ we decompose
$$
U_i \cap U_j = \coprod U_{i, j, l}
$$
into its finitely many connected components. Write
$K_{i, j, l} \subset \mathcal{O}(U_{i, j, l})$
for the integral closure of $k_1$ and of $k_2$. By
Lemma \href{https://stacks.math.columbia.edu/tag/04MI}{lemma integral closure ground field (Tag 04MI)}
the rings $K_i$ and $K_{i, j, l}$ are fields.
Now we claim that $k_1'$ and $k_2'$ both equal the kernel of the map
$$
\prod K_i \longrightarrow \prod K_{i, j, l}, \quad
(x_i)_i \longmapsto x_i|_{U_{i, j, l}} - x_j|_{U_{i, j, l}}
$$
which proves what we want.
Namely, it is clear that $k_1'$ is contained in this kernel.
On the other hand, suppose that $(x_i)_i$ is in the kernel.
By the sheaf condition $(x_i)_i$ corresponds to $f \in \mathcal{O}(X)$.
Pick some $i_0 \in I$ and let $P(T) \in k_1[T]$ be a monic polynomial
with $P(x_{i_0}) = 0$. Then we claim that $P(f) = 0$ which proves
that $f \in k_1$. To prove this we have to show that $P(x_i) = 0$
for all $i$. Pick $i \in I$. As $X$ is connected there exists a
sequence $i_0, i_1, \ldots, i_n = i \in I$ such that
$U_{i_t} \cap U_{i_{t + 1}} \not = \emptyset$. Now this means that
for each $t$ there exists an $l_t$ such that $x_{i_t}$ and
$x_{i_{t + 1}}$ map to the same element of the field $K_{i_t, i_{t + 1}, l_t}$.
Hence if $P(x_{i_t}) = 0$, then $P(x_{i_{t + 1}}) = 0$. By
induction, starting with $P(x_{i_0}) = 0$ we deduce that
$P(x_i) = 0$ as desired.

\medskip\noindent
To finish the proof of the lemma we prove the lemma under the
additional hypothesis that $X$ is quasi-compact. By
Lemma \href{https://stacks.math.columbia.edu/tag/04MI}{lemma integral closure ground field (Tag 04MI)}
after replacing $k_i$ by $k_i'$
we may assume that $k_i$ is integrally closed in $\Gamma(X, \mathcal{O}_X)$.
This implies that $\mathcal{O}(X)^*/k_i^*$ is a finitely generated
abelian group, see
Proposition \href{https://stacks.math.columbia.edu/tag/04L7}{proposition units general (Tag 04L7)}.
Let $k_{12} = k_1 \cap k_2$ as a subring of $\mathcal{O}(X)$.
Note that $k_{12}$ is a field. Since
$$
k_1^*/k_{12}^* \longrightarrow \mathcal{O}(X)^*/k_2^*
$$
we see that $k_1^*/k_{12}^*$ is a finitely generated abelian group as well.
Hence there exist $\alpha_1, \ldots, \alpha_n \in k_1^*$ such that
every element $\lambda \in k_1$ has the form
$$
\lambda = c \alpha_1^{e_1} \ldots \alpha_n^{e_n}
$$
for some $e_i \in \mathbf{Z}$ and $c \in k_{12}$.
In particular, the ring map
$$
k_{12}[x_1, \ldots, x_n, \frac{1}{x_1 \ldots x_n}] \longrightarrow k_1, \quad
x_i \longmapsto \alpha_i
$$
is surjective. By the Hilbert Nullstellensatz,
Algebra, Theorem \href{https://stacks.math.columbia.edu/tag/00FV}{theorem nullstellensatz (Tag 00FV)}
we conclude that $k_1$ is a finite extension of $k_{12}$.
In the same way we conclude that $k_2$ is a finite extension of $k_{12}$.
In particular both $k_1$ and $k_2$ are contained in the integral closure
$k_{12}'$ of $k_{12}$ in $\Gamma(X, \mathcal{O}_X)$. But since $k_{12}'$
is a field by
Lemma \href{https://stacks.math.columbia.edu/tag/04MI}{lemma integral closure ground field (Tag 04MI)}
and since we chose $k_i$ to be integrally closed in $\Gamma(X, \mathcal{O}_X)$
we conclude that $k_1 = k_{12} = k_2$ as desired.
\end{proof}
\end{document}
